\subsection{Технологии поиска и обработки естественного языка в системах RAG}

Использование генеративных моделей в чистом виде
влечет высокие риски «конфабуляций» и «галлюцинаций»:
ответ модели может казаться правдоподобным, но содержать фактологические
или причинно-следственные ошибки~\cite{Ji_2023}.

Для решения этой проблемы была разработана технология
Retrieval-Augmented Generation (RAG) \cite{Lewis_2020}.
Она добавляет дополнительный этап поиска информации перед генерацией ответа.
На основе запроса пользователя алгоритм извлекает наиболее релевантную
информацию из доступных источников, таких как текстовые файлы или базы данных,
после чего передает ее языковой модели в качестве контекста.
Такой подход заставляет LLM выступать в роли анализатора,
а не источника фактов.

Эффективность RAG-систем зависит от качества алгоритмов извлечения.
В связи с этим, существует несколько подходов к поиску необходимых данных.
Наиболее популярными методами являются лексический и семантический поиск.

Лексический поиск основан на точных совпадениях
лексем из запроса пользователя с содержащимися в хранилищах.
Данный подход отлично справляется с поиском специфических терминов,
имен собственных и прочих точных формулировок, но
не справляется, например, в случае замены слов на синонимы или чередующихся корней.
В случае нахождения лексемы в корпусе из него извлекается
фрагмент из некоторого количества предшествующих и последующих слов или предложений.
Популярным алгоритмом лексического поиска является BM25~\cite{Robertson_2009}.

Для преодоления этих ограничений используется семантический поиск,
основанный на векторном сходстве~\cite{Karpukhin_2020}.
Исходный текст разбивается на фрагменты (чанки), для каждого из которых вычисляется
многомерный числовой вектор (эмбеддинг) из одного пространства.
Поиск происходит путем вычисления сходства между этими векторами
с помощью метрик, таких как косинусное сходство.
Если полученное значение больше некоторого порога, то
чанки считаются релевантными, и текст, соответствующих эмбеддингу,
извлекается.
Этот подход позволяет улавливать смысл текста,
не основываясь на конкретных словах или формулировках.

Векторный поиск также имеет несколько недостатков~\cite{Gao_2023}:
\begin{itemize}
      \item Потеря конкретики: смысловое значение фрагментов усредняется,
            что ведет к потере деталей, таких как имен собственных, использованных терминов, и так далее.
            Чем больше чанк, тем больше деталей теряется.
      \item Проблема «дальних эмбеддингов»: схожие данные могут иметь сильно
            отличающиеся друг от друга векторы, из-за чего могут не быть найдены.
\end{itemize}

Современные RAG системы комбинируют эти подходы в
так называемый гибридный поиск~\cite{Gao_2023}. Он позволяет
компенсировать недостатки одного подхода с помощью другого,
что потенциально может увеличить полноту поиска.
В нем извлеченные чанки получают унифицированную оценку
и ранжируются по ней.

Несмотря на все преимущества этих подходов,
обработка сложных аналитических запросов,
требующих многошагового рассуждения,
требует более сложных алгоритмов поиска.
В случаях, когда нужная информация напрямую не связана с запросом,
она не может быть найдена
ни семантическим, ни лексическим поиском, из-за чего
поисковый алгоритм RAG не сможет извлечь необходимые данные.

Более продвинутой стратегией поиска является Graph RAG~\cite{Edge_2024}.
В ней в алгоритм поиска дополняется обходом графа знаний.
Для первоначального сбора вершин используется гибридный поиск.
Затем начинается итеративный обход соседей на основе
имеющихся ребер и их типов, что позволяет системе собирать дополнительный
контекст из файлов, которые могут не содержать ключевых слов, или быть семантически далекими,
но при этом логически связанными и обладающими необходимой информацией,
что позволяет значительно расширить контекст.

\textbf{Языковые модели в задаче связывания.}
Поскольку необходимо искать тексты, обладающие не только высоким семантическим сходством
или содержащие ключевые слова, нужно использовать иные от сравнения эмбеддингов подходы для анализа
связности текстов. Для решения этой задачи используются модели
машинного обучения, такие как Cross Encoder~\cite{Nogueira_2019} и большие языковые модели (LLM).
Они позволяют реранжировать тексты на основе оценки наличия причинно-следственных, логических
и иных сложных связей.

Cross Encoder являются сравнительно небольшими моделями для попарной проверки текстов.
Эти модели принимают на вход два текста и совместно их обрабатывают с помощью cross-attention,
что позволяет выявить более тонкие логические связи и точнее провести фильтрацию.

Большие языковые модели
отлично понимают естественный язык и способны выявлять нетривиальные, многошаговые
логические цепочки, которые могли быть упущены при фильтрации через Cross Encoder.
Кроме того, LLM могут комбинировать фильтрацию с типизацией отношений,
что позволяет решить обе задачи за один запрос.
Однако такой подход является вычислительно сложным и, как следствие, дорогим, из-за чего
LLM используют только для классификации и дополнительной фильтрации уже отфильтрованных пар.
